\chapter{Evaluating agents honestly}
\label{ch:eval}

An impressive demo proves nothing. This chapter is the measurement discipline that makes the
project's claims trustworthy, drawn from the agent-evaluation literature.

\section{Reliability: pass@k vs pass\textasciicircum k}

\term{pass@k} asks whether the agent succeeds in \emph{at least one} of $k$ tries; it flatters
flaky agents. \term{pass\^{}k} asks whether it succeeds in \emph{all} $k$ tries; it measures
reliability~\citep{tau_bench}. As $k$ grows the two diverge sharply for an unreliable agent.

\begin{figure}[h]\centering
\begin{tikzpicture}
\begin{axis}[width=9.5cm,height=5.4cm,xlabel={$k$ (independent attempts)},ylabel={success rate},
  xmin=1,xmax=8,ymin=0,ymax=1.02,legend pos=south west,legend cell align=left,
  xtick={1,2,3,4,5,6,7,8}, grid=both, grid style={clExternal!20}]
  % per-attempt success p=0.7
  \addplot[very thick, clBackend, mark=*, mark size=1] coordinates
    {(1,0.7)(2,0.91)(3,0.973)(4,0.992)(5,0.998)(6,0.999)(7,1.0)(8,1.0)};
  \addlegendentry{pass@k (at least one)}
  \addplot[very thick, clSecurity, mark=square*, mark size=1] coordinates
    {(1,0.7)(2,0.49)(3,0.343)(4,0.240)(5,0.168)(6,0.118)(7,0.082)(8,0.058)};
  \addlegendentry{pass\^{}k (all $k$)}
\end{axis}
\end{tikzpicture}
\caption{For a per-attempt success of 0.7, pass@k rushes to 1 while pass\textasciicircum k
collapses. Reliability is the honest metric.\datalabel{Exact (computed from $p=0.7$)}}
\label{fig:passk}
\end{figure}

\section{Cost per successful task and the Pareto frontier}

Accuracy alone hides the bill. The right object is the \term{accuracy--cost Pareto frontier}:
plot each system by cost and success, and compare frontiers, not single
points~\citep{ai_agents_that_matter}. A famous result: a plain ``retry'' baseline beat an
elaborate agent at a fraction of the cost, which is why baselines like ARM-A5 (random) and
ARM-A6 (matched compute) exist in our design.

\begin{figure}[h]\centering
\begin{tikzpicture}
\begin{axis}[width=9.5cm,height=5.6cm,xlabel={cost per task (USD, log)},ylabel={pass\^{}4},
  xmode=log, xmin=0.002,xmax=1,ymin=0.3,ymax=0.9,legend pos=south east,
  legend cell align=left, grid=both, grid style={clExternal!20},
  nodes near coords, point meta=explicit symbolic, every node near coord/.append style={font=\tiny}]
  \addplot[only marks, mark=*, clModel, mark size=1.6] coordinates
    {(0.004,0.62)[A1 hybrid] (0.06,0.71)[A0 LLM-only] (0.006,0.60)[A2 same-iface]};
  \addplot[only marks, mark=square*, clExternal, mark size=1.6] coordinates
    {(0.09,0.72)[A6 matched] (0.004,0.40)[A5 random]};
\end{axis}
\end{tikzpicture}
\caption{Illustrative accuracy--cost plot. The question is whether the hybrid (A1) sits
up-and-left of the baseline (A0) and beats the random control (A5).\datalabel{Illustrative (synthetic data); arm definitions from docs/PLAN.md \S9}}
\label{fig:pareto}
\end{figure}

\section{The controls that stop us fooling ourselves}

\begin{description}[leftmargin=1.4cm,style=nextline]
  \item[Same-interface baseline (ARM-A2)] an LLM answering the \emph{same} typed questions,
    to separate ``typed decision points help'' from ``Jev specifically helps.''
  \item[Matched compute (ARM-A6)] give the baseline the same extra budget the hybrid spends,
    so a win is not just ``more compute.''
  \item[Random policy (ARM-A5)] the same machinery making random decisions, to control for
    the structure of the intervention.
  \item[Oracle (ARM-A4)] perfect decisions, an upper bound: if even the oracle does not help,
    the chosen decision points are not the bottleneck.
  \item[Pre-registration] fix arms, metrics, thresholds, and seeds \emph{before} looking at
    results; fit thresholds on a disjoint split; report bootstrap confidence intervals over
    $\ge 3$ seeds.
\end{description}

\begin{pitfall}
Never use an LLM judge as sole ground truth for accuracy: it measures \emph{agreement}, not
correctness. One study reported 91.5\% agreement with a frontier model, but with no human
labels that is agreement, not accuracy~\citep{good_start_labs}.
\end{pitfall}

\begin{checkbox}
An agent reports pass@8 of 0.95. Why might its pass\^{}8 be under 0.6, and which figure shows
it? (Appendix~\ref{app:answers}.)
\end{checkbox}
